\headline{{\textbf{\Large\sffamily Seasonal Care Tips:}}\\
          {\textbf{\large\sffamily Mid\--October to Mid\--November}}}
\label{2025-10-care}
\noindent
From mid\--October through mid\--November, autumn deepens into its reflective phase.
The first storm of the season has swept through, bringing cold air, more rain, and forceful winds that have hastened leaf drop.
Trees stand increasingly bare, the year's growth complete, their energies drawn inward.
Our focus now turns from refinement to preservation---protecting roots, structure, and soil as temperatures fall and daylight wanes.
Calm observation and small, consistent actions will ensure healthy dormancy and a strong start next spring.

\medskip
\topic{Wind, Rain, and Structural Stability}
\noindent
The recent storm serves as a reminder of autumn's volatility.
Secure trees that rock in their pots, and check for root exposure or tilted soil.
Use light ties or bench anchors for shallow containers prone to toppling.

If bark has been scraped or branches torn, trim cleanly and seal wounds to prevent infection.
Inspect roots near pot edges---if fine roots are visible, cover gently with fresh substrate rather than compressing wet soil.
Persistent rain can compact mixes and reduce airflow; a gentle stir with a chopstick helps restore porosity without disturbing roots.

\medskip
\topic{Watering and Airflow in a Damp Month}
\noindent
With regular rain, overwatering is now a greater threat than drought.
Check soil daily but water only when the surface begins to dry or lighten in colour.
Sheltered trees under eaves or polycarbonate roofs may dry more quickly---judge each individually.

Raise pots slightly off benches to aid drainage and prevent contact with standing water.
Avoid trays or saucers, which trap cold moisture beneath containers.
Conifers, junipers, and olives prefer a drier footing and should be protected from persistent rain, ideally under transparent covers that still admit light.

\medskip
\topic{Fertilising: Winding Down for Dormancy}
\noindent
By now, deciduous trees have little appetite for nutrients.
Cease all feeding once leaves have fallen, and remove any remaining fertiliser cakes to avoid mould.

Evergreens, however, still benefit from one final, weak application of a low\--nitrogen, potassium\--rich feed.
This helps strengthen roots and sustain colour through winter.
After that, allow soil chemistry to stabilise naturally.
Do not fertilise again until buds begin to swell in spring; excess nutrients now only sour the substrate.

\medskip
\topic{Fungal Problems and Pests in Cool, Wet Weather}
\noindent
Moist, decaying debris invites fungus and pests.
Sweep leaves regularly from benches and pot surfaces, especially after wind or rain.
Do not compost diseased foliage---dispose of it away from the collection.

Watch for powdery mildew on maple, rust on hawthorn, and mould developing around fertiliser remnants.
Scrape off any white growth and top up with clean, dry akadama or gravel.
Vine weevil larvae, slugs, and root aphids may still be active; inspect soil edges and root collars closely.
If mild weather persists, biological controls remain effective for a few more weeks.

\medskip
\topic{Pruning, Wiring, and Clean\--Up Work}
\noindent
With most deciduous trees bare, their framework is clearly visible.
Use this time for quiet inspection and light maintenance rather than structural change.
Remove dead twigs or storm damage, but leave shaping cuts for winter dormancy or early spring.

Wiring can still be done on healthy material, though branches are brittle in cool conditions.
Apply gently, avoiding sharp bends, and use aluminium wire where possible---it resists corrosion better than copper in damp air.
After wiring, store trees under partial shelter to keep wires clean and rust free.
This is also a good moment to clean tools, benches, and pots before winter storage.

\medskip
\topic{Positioning and Cold Protection}
\noindent
Light levels diminish quickly now, but exposure remains important.
Place trees where they receive full morning sun to dry foliage and discourage mildew.
After each storm, realign pots and resecure loose coverings.

Begin separating species by hardiness.
Hardy natives can remain outside with wind protection, while tender or Mediterranean species should move into cold frames or unheated greenhouses as nights cool below 8–10°C.
Wrap vulnerable pots with fleece or bubble wrap to buffer roots against rapid temperature drops.
Maintain airflow---stagnant, humid air is far more damaging than cold.

\medskip
\topic{Leaf Drop and Dormancy Observation}
\noindent
Rapid yellowing and leaf fall are natural responses to the season's chill.
Allow the process to unfold without interference; remaining leaves continue to feed roots until they detach.

Once bare, examine branches for insect eggs and wipe them away with a soft brush.
Note which trees entered dormancy first---this helps refine watering and repotting schedules next year.
As metabolism slows, reduce watering further but never let soil dry completely.
Cool, slightly moist soil encourages stable dormancy and healthy root function.

\medskip
\topic{Preparing for Winter}
\noindent
By early November, winter routines should be nearly in place.
Check cold frames, fleece wraps, and rain shelters for damage after each windy spell.
Ensure that lids and vents operate smoothly---ventilate on mild days to prevent condensation.

Clean under benches and walkways to discourage pests and algae.
Add a layer of coarse gravel beneath benches to improve drainage.
Gather fleece, mesh, and frost cloth in readiness for sudden cold snaps.
This quieter interval, after the year's activity but before deep winter, rewards unhurried care and observation.
Take time to appreciate the architecture of bare branches and the subdued beauty of the resting garden.
Each tree, stripped of adornment, holds the promise of renewal within its silent form.
\closearticle
%\columnbreak
